\documentclass[12pt, a4paper]{article}

% ─── Packages ────────────────────────────────────────────────────────────────
\usepackage[a4paper, top=2.5cm, bottom=2.5cm, left=2.5cm, right=2.5cm]{geometry}
\usepackage{amsmath, amssymb, amsfonts}
\usepackage{bm}
\usepackage{graphicx}
\usepackage{xcolor}
\usepackage{booktabs}
\usepackage{array}
\usepackage{multirow}
\usepackage{hyperref}
\usepackage{caption}
\usepackage{subcaption}
\usepackage{setspace}
\usepackage{parskip}
\usepackage{enumitem}
\usepackage{titlesec}
\usepackage{fancyhdr}
\usepackage{abstract}
\usepackage[numbers, sort&compress]{natbib}
\usepackage{microtype}
\usepackage{float}
\usepackage[section]{placeins}

% TikZ
\usepackage{tikz}
\usepackage{pgfplots}
\pgfplotsset{compat=1.18}
\usetikzlibrary{
	shapes.geometric, shapes.misc, shapes.arrows,
	arrows.meta, positioning, fit, calc,
	backgrounds, decorations.pathreplacing,
	matrix, chains, shadows.blur
}

% ─── Colours ─────────────────────────────────────────────────────────────────
\definecolor{darkblue}{RGB}{31,56,100}
\definecolor{midblue}{RGB}{46,95,163}
\definecolor{lightblue}{RGB}{173,210,240}
\definecolor{accentgreen}{RGB}{0,128,96}
\definecolor{accentorange}{RGB}{210,100,20}
\definecolor{accentpurple}{RGB}{120,60,160}
\definecolor{lightgray}{RGB}{245,245,245}
\definecolor{midgray}{RGB}{180,180,180}
\definecolor{blockA}{RGB}{214,234,248}
\definecolor{blockB}{RGB}{212,239,223}
\definecolor{blockC}{RGB}{252,235,212}
\definecolor{blockD}{RGB}{235,218,252}
\definecolor{blockE}{RGB}{252,212,212}

% ─── Hyperref ────────────────────────────────────────────────────────────────
\hypersetup{
	colorlinks=true, linkcolor=midblue, citecolor=accentgreen,
	urlcolor=midblue, pdftitle={Multi-Method EMG Prediction from EEG on WAY-EEG-GAL}
}

% ─── Typography ──────────────────────────────────────────────────────────────
\setstretch{1.25}
\titleformat{\section}{\large\bfseries\color{darkblue}}{\thesection}{1em}{}
\titleformat{\subsection}{\normalsize\bfseries\color{midblue}}{\thesubsection}{1em}{}
\titleformat{\subsubsection}{\normalsize\bfseries\color{midblue}}{\thesubsubsection}{1em}{}

\setlength{\headheight}{13.6pt}
\addtolength{\topmargin}{-1.6pt}
\pagestyle{fancy}
\fancyhf{}
\rhead{\small\textcolor{midgray}{Multi-Method EEG-to-EMG Prediction on WAY-EEG-GAL}}
\lhead{\small\textcolor{midgray}{\thepage}}
\renewcommand{\headrulewidth}{0.4pt}

% ─── Math operators ──────────────────────────────────────────────────────────
\DeclareMathOperator{\Cov}{Cov}
\DeclareMathOperator{\Var}{Var}
\DeclareMathOperator{\MSE}{MSE}
\DeclareMathOperator{\softmax}{softmax}
\newcommand{\R}{\mathbb{R}}

% ─────────────────────────────────────────────────────────────────────────────
\begin{document}
	
	% ─── TITLE ───────────────────────────────────────────────────────────────────
	\begin{center}
		{\LARGE\bfseries\color{darkblue}
			Predicting EMG Envelopes from EEG and Kinematic Signals:}\\[14pt]
		{\large\itshape\color{midblue}
			A Research Article Proposal}\\[20pt]
		{\normalsize
			\textbf{\href{mailto:sharifbeigymohammad@gmail.com}{MohammadMahdi SharifBeigi}}\\[4pt]
			\textit{Computational Neuroscience Team 6}}\\[16pt]
	\end{center}
	
	\vspace{6pt}
	\hrule height 1pt
	\vspace{14pt}
	
	% ─── ABSTRACT ────────────────────────────────────────────────────────────────
	\begin{abstract}
		\noindent
		Electromyography (EMG)-based prosthetic control fails in cases of
		high-level amputation or severe paralysis where functional muscles are
		absent. Brain-derived signals from electroencephalography (EEG) offer
		a complementary pathway. This article presents a \textbf{multi-method
		comparative study} for continuous prediction of EMG envelopes from
		scalp EEG and kinematic measurements during grasp-and-lift tasks,
		validated on the \textbf{WAY-EEG-GAL} dataset (3,936 trials, 12
		participants, 32-channel EEG, 5-channel EMG, 3D kinematics, 6-axis
		force/torque). We propose and compare four architectures of
		increasing complexity:
		
		\textbf{Method~1 (KG-GT):} A Kinematic-Guided Graph-Transformer
		combining CCA-based EEG--kinematic alignment, a Transformer encoder,
		and a kinematic-conditioned GAT with a hybrid MSE~+~SoftDTW loss.
		
		\textbf{Method~2 (NNMF-KG-GT):} Enhances Method~1 by replacing the
		anatomically fixed graph structure with Non-Negative Matrix
		Factorisation (NNMF)-derived muscle synergy weights, providing
		biologically interpretable, data-driven inter-muscle connectivity.
		
		\textbf{Method~3 (PTL-Net):} A Phase-Targeted Lightweight Network
		that focuses on Load and Hold movement phases, uses correlation-based
		channel selection (6 channels: C3, C4, Cz, CP3, CP4, Fz), and
		performs multi-task prediction of EMG envelopes, Grip Force, and
		object weight classification---demonstrating direct clinical
		applicability for adaptive prosthetic control.
		
		\textbf{Method~4 (USPT):} A Unified Synergy-Phase-Transformer that
		integrates all innovations from Methods~1--3 into a single framework
		with NNMF-initialised kinematic-guided GAT, phase-aware Transformer
		encoding, triple-head decoding, and phase-weighted loss.
		
		\vspace{4pt}
		\noindent\textbf{Keywords:} EEG, EMG prediction, WAY-EEG-GAL,
		NNMF, muscle synergy, graph attention network, Transformer,
		phase-targeted decoding, multi-task learning, prosthetic control.
	\end{abstract}
	
	\hrule height 0.5pt
	\vspace{10pt}
	
	\tableofcontents
	\newpage
	
	% ═════════════════════════════════════════════════════════════════════════════
	\section{Introduction}
	% ═════════════════════════════════════════════════════════════════════════════
	
	The ability to decode motor intent from neural signals and translate it into
	reliable control of external devices is a foundational challenge in neural
	engineering. Two principal peripheral modalities have historically dominated
	prosthetic control: surface electromyography (sEMG), which measures the
	electrical activity of muscles, and joint kinematics, which captures the
	positional state of limbs. While both are effective for able-bodied
	individuals and low-level amputees, their availability becomes severely
	restricted in cases of high-level limb loss, spinal cord injury, or
	progressive neuromuscular disease.
	
	Electroencephalography (EEG) offers a non-invasive, low-cost alternative by
	capturing oscillatory cortical activity through scalp-mounted electrodes.
	Motor-related cortical areas exhibit characteristic modulations in the
	mu~($8$--$13$~Hz) and beta~($13$--$30$~Hz) frequency bands during movement
	preparation and execution, as well as slow cortical potentials in the delta
	band ($0.1$--$3$~Hz) that track the continuous dynamics of movement.
	
	The core scientific question addressed here is: \textit{to what extent can
		the continuous EMG envelope be accurately predicted from scalp EEG
		recordings, and can the simultaneous inclusion of kinematic measurements,
		muscle synergy decomposition, and phase-targeted processing further
		improve this prediction?}
	
	This work makes the following contributions:
	\begin{itemize}[itemsep=2pt]
		\item A novel \textbf{CCA-based EEG--kinematic alignment} stage (Method~1) and a \textbf{Corticomuscular Coherence (CMC)-based EEG--EMG connectivity} feature extraction (Method~4) that provide a biologically-grounded projection before deep learning.
		\item A \textbf{time-lagged windowing strategy} to account for physiological delays between cortical planning, muscle activation, and limb kinematics~\cite{esi_gal_2025}.
		\item A \textbf{Transformer encoder} that explicitly models the nonlinear,
		multi-lag temporal coupling between cortical activity and muscle
		activation (Methods~1, 2, 4).
		\item A \textbf{Kinematic-Guided GAT} whose attention coefficients are
		dynamically conditioned on the instantaneous kinematic state, enabling
		task-phase-dependent muscle synergy modelling (Methods~1, 2, 4).
		\item Integration of \textbf{NNMF-derived muscle synergy} weights as
		data-driven graph adjacency priors, inspired by multimodal approaches
		to motor dynamics~\cite{pierella2020} (Methods~2, 4).
		\item A \textbf{phase-targeted, channel-selected, multi-task} architecture
		that predicts EMG envelopes, Grip Force, and object weight class
		simultaneously from only 6 EEG channels during Load and Hold phases
		(Methods~3, 4).
		\item A \textbf{Unified Synergy-Phase-Transformer (USPT)} combining
		all innovations with phase-weighted loss and staged training (Method~4).
		\item A \textbf{hybrid MSE~+~SoftDTW loss} function that penalises both
		amplitude error and temporal misalignment (all methods).
		\item Comprehensive validation on the richly multimodal
		\textbf{WAY-EEG-GAL} dataset across 12 subjects and systematically
		varied grasp conditions.
	\end{itemize}
	
	% ═════════════════════════════════════════════════════════════════════════════
	\section{Background and Related Work}
	% ═════════════════════════════════════════════════════════════════════════════
	
	\subsection{The WAY-EEG-GAL Dataset}
	
	The WAY-EEG-GAL dataset~\cite{luciw2014} contains 3,936 grasp-and-lift (GAL)
	trials from 12 right-handed participants. The participant grasped a small
	instrumented object with the thumb and index finger, lifted it approximately
	5~cm, held it at a target position, and replaced it. Object weight was varied
	between 165, 330, and 660~g; contact surface was varied between sandpaper,
	suede, and silk---in an unpredictable manner---enforcing systematic changes
	in fingertip force coordination. Simultaneous recordings include:
	\begin{itemize}[itemsep=2pt]
		\item \textbf{EEG:} 32 channels (ActiCap), sampled at 500~Hz.
		\item \textbf{EMG:} 5 channels (anterior deltoid, brachioradialis, flexor
		digitorum, extensor digitorum, first dorsal interosseus), sampled at
		4,000~Hz.
		\item \textbf{Kinematics:} Four Polhemus FASTRAK 3D position/orientation
		sensors (object, index finger, thumb, wrist), sampled at 500~Hz.
		\item \textbf{Force/torque:} Two ATI Nano-17 6-axis sensors at both contact
		plates.
	\end{itemize}
	Sixteen behaviourally relevant events were extracted per trial, and all
	participants were confirmed to adapt grip coordination to the prevailing
	weight and friction~\cite{luciw2014}.
	
	\subsection{Unscented Kalman Filtering Approaches}
	
	Nakagome et al.~\cite{nakagome2017} demonstrated offline continuous decoding of
	lower-limb EMG envelopes during over-ground walking on multiple terrains from
	delta-band ($0.1$--$3$~Hz) EEG using a first-order Unscented Kalman
	Filter~(UKF).  Across six subjects the mean Pearson correlation was
	$r = 0.236$, with $79.3\%$ of results exceeding $\mathrm{SNR} = 0$~dB.
	
	Sburlea, Butturini, and M\"{u}ller-Putz~\cite{sburlea2021} extended the UKF
	approach to upper-limb grasping (31 subjects, 33 grasping conditions) using
	low-frequency EEG ($0.1$--$2$~Hz), a multivariate autoregressive (MVAR) state
	model, and partial least squares (PLS) neural tuning, achieving median
	$r = 0.20$ and $\mathrm{nMAE} = 28.6\%$.
	
	The UKF state-transition equation is:
	\begin{equation}
		\bm{x}_n = \sum_{i=1}^{m} \bm{A}_i(n)\,\bm{x}_{n-i} + \bm{e}_n,
		\label{eq:ukf_state}
	\end{equation}
	where $\bm{x}_n \in \R^d$ is the EMG envelope state, $\bm{A}_i(n)$ are
	$d\!\times\!d$ coefficient matrices, and $\bm{e}_n$ is Gaussian noise.
	The observation model is:
	\begin{equation}
		\bm{y}_k = \bm{H}\,\bm{x}_k + \bm{v}_{k-1},
		\label{eq:ukf_obs}
	\end{equation}
	where $\bm{H}$ is the PLS-estimated neural tuning matrix and $\bm{v}$ is
	measurement noise. The UKF propagates sigma points through
	Eq.~\eqref{eq:ukf_state} to estimate state statistics nonlinearly.
	However, the \emph{fixed linear} observation model $\bm{H}$ and the
	Gaussian noise assumptions limit performance ceiling for complex grasping.
	
	\subsection{Deep Sequence Learning: CNN-LSTM}
	
	Jain et al.~\cite{jain2025} proposed \textit{RawSeq2SeqNet}, a CNN-LSTM
	architecture for reconstructing biceps brachii EMG from 19-channel EEG during
	elbow flexion-extension (single subject). The key innovation was a hybrid
	training loss:
	\begin{equation}
		\mathcal{L} = 0.9 \cdot \MSE + 0.1 \cdot \mathrm{SoftDTW}(\gamma=0.2),
		\label{eq:jain_loss}
	\end{equation}
	achieving $R^2 = 0.364$ and Pearson $r = 0.795$ on the held-out test set, with
	mean inference time $2.11$~ms per 2-second window. The SoftDTW component
	tolerates the physiological cortico-muscular conduction delay
	($\approx 10$--$30$~ms) during training.
	
	\subsection{Graph-Based Architectures with Muscle Synergy}
	
	Li et al.~\cite{li2025} introduced \textit{MSGAT-LSTM} for continuous wrist
	kinematics prediction from sparse sEMG, using a GAT with edge weights
	derived from cosine similarity of EMG feature vectors, grounded in muscle
	synergy (MS) theory. The cosine similarity between nodes $i$ and $j$ after
	shared linear transform $\bm{W}$ is:
	\begin{equation}
		\alpha_{ij} = \frac{\sum_{k=1}^{F'} m_{ik}\,m_{jk}}
		{\|\bm{M}_i\|\,\|\bm{M}_j\|},
		\label{eq:cosine}
	\end{equation}
	with normalised edge weights:
	\begin{equation}
		e_{ij} = \frac{\exp(\mathrm{LeakyReLU}(\alpha_{ij}))}
		{\sum_{j'} \exp(\mathrm{LeakyReLU}(\alpha_{ij'}))}.
		\label{eq:gat_edge}
	\end{equation}
	On NinaPro DB2 (12 subjects), MSGAT-LSTM achieved $R^2 = 0.960$ and
	$\mathrm{RMSE} = 0.098$~rad using only 10 sparse electrodes, reducing
	training time by $\approx 13\%$ relative to GCN-LSTM.
	
	\subsection{Non-Negative Matrix Factorisation for Muscle Synergy}
	
	Pierella et~al.~\cite{pierella2020} demonstrated a multimodal approach to
	capturing post-stroke motor recovery dynamics by integrating EEG, EMG, and
	kinematic data. A central analytical tool was \textbf{Non-Negative Matrix
	Factorisation (NNMF)}~\cite{lee1999}, which decomposes the multi-muscle EMG
	activation matrix $\bm{V} \in \R_{\geq 0}^{T \times M}$ into a small
	number $r$ of synergy modules:
	\begin{equation}
		\bm{V} \approx \bm{W}\,\bm{H},
		\qquad \bm{W} \in \R_{\geq 0}^{T \times r},\;
		\bm{H} \in \R_{\geq 0}^{r \times M},
		\label{eq:bg_nnmf}
	\end{equation}
	where $\bm{W}$ encodes temporal activation profiles and $\bm{H}$ encodes
	muscle recruitment weights for each synergy. The non-negativity constraint
	ensures physiological interpretability: muscles can only be activated, not
	inhibited, by a synergy module.
	
	Pierella et~al.\ showed that the number of synergies, their composition,
	and their temporal activation patterns change during post-stroke
	rehabilitation, tracking recovery of motor coordination at the spinal level.
	Tresch et~al.~\cite{tresch2006} provided foundational evidence that the
	central nervous system controls movement through a small set of
	muscle synergies, validating NNMF as a principled decomposition for
	motor neuroscience.
	
	\paragraph{Gap.}
	Despite the established utility of NNMF for synergy analysis, no prior
	EEG-to-EMG prediction study has used NNMF-derived synergy weights as
	structural priors for graph-based neural network architectures. This gap
	is addressed by Methods~2 and~4 in the present work.
	
	% ═════════════════════════════════════════════════════════════════════════════
	\section{Research Gaps and Motivation}
	% ═════════════════════════════════════════════════════════════════════════════
	
	The literature review reveals eight specific gaps:
	
	\begin{enumerate}[label=\textbf{G\arabic*.}, leftmargin=*, itemsep=4pt]
		\item \textbf{Limited expressivity.} UKF imposes linear, Gaussian
		assumptions on the EEG--EMG mapping; deep architectures with
		transformer-level capacity are needed.
		\item \textbf{Single-subject scope.} RawSeq2SeqNet was validated on one
		participant; the WAY-EEG-GAL dataset provides 12 participants and 3,936
		trials.
		\item \textbf{Underutilisation of kinematics.} No existing EEG-to-EMG
		predictor leverages the rich kinematic data (3D finger positions, grip
		aperture, load/grip force) available in WAY-EEG-GAL as a structural
		constraint.
		\item \textbf{No synergy structure in EEG decoding.} The muscle-synergy
		graph principle proven in MSGAT-LSTM has never been applied to the
		EEG-to-EMG direction.
		\item \textbf{Temporal misalignment.} With the exception of Jain et al.,
		prior EEG-to-EMG methods use MSE/MAE loss functions sensitive to
		cortico-muscular conduction delays.
		\item \textbf{No NNMF-derived priors for graph construction.} Existing
		graph-based architectures use either anatomically fixed or learned-from-scratch
		graph topologies. NNMF-derived synergy weights---validated by decades of
		motor neuroscience research~\cite{tresch2006, pierella2020}---have
		never been used to initialise the adjacency matrix of a GAT for
		EEG-to-EMG prediction.
		\item \textbf{Full-window prediction ignores phase-specific coupling.}
		All prior methods predict EMG uniformly across the entire trial window.
		However, the EEG--EMG relationship is strongest during Load and Hold
		phases, where sustained cortical drive maintains isometric contraction.
		Reach and Replace phases have weaker coupling and higher artefact
		contamination.
		\item \textbf{No multi-task or classification from EEG.} No EEG-to-EMG
		study simultaneously predicts EMG envelopes, grip force, and object
		weight class. Multi-task learning can improve representation quality
		through shared features, and weight classification directly demonstrates
		clinical applicability for adaptive prosthetic grip control.
	\end{enumerate}
	
	\newpage
	
	% ═════════════════════════════════════════════════════════════════════════════
	\section{Methods Overview}
	% ═════════════════════════════════════════════════════════════════════════════
	
	This work proposes four methods of increasing complexity, each addressing
	a progressively larger subset of the identified research gaps. Table~\ref{tab:methods_overview}
	summarises the four methods.
	
	\begin{table}[H]
		\centering
		\caption{Overview of the four proposed methods.}
		\label{tab:methods_overview}
		\renewcommand{\arraystretch}{1.4}
		\begin{tabular}{clp{5.5cm}c}
			\toprule
			\textbf{Method} & \textbf{Name} & \textbf{Key Innovation} & \textbf{Gaps Addressed} \\
			\midrule
			1 & KG-GT &
			CCA + Transformer + Kinematic-Guided GAT + SoftDTW &
			G1--G5 \\
			2 & NNMF-KG-GT &
			Method~1 + NNMF-derived synergy graph adjacency &
			G1--G6 \\
			3 & PTL-Net &
			Phase targeting + channel selection + multi-task (EMG + Force + weight class) &
			G7--G8 \\
			4 & USPT &
			Unified: NNMF + phase embeddings + channel selection + CMC + KG-GAT + triple-head &
			G1--G9 \\
			\bottomrule
		\end{tabular}
	\end{table}
	
	% ═════════════════════════════════════════════════════════════════════════════
	\section{Method~1: Kinematic-Guided Graph-Transformer (KG-GT)}
	\label{sec:method1}
	% ═════════════════════════════════════════════════════════════════════════════
	
	\input{method1}
	
	\newpage
	
	% ═════════════════════════════════════════════════════════════════════════════
	\section{Method~2: NNMF-Enhanced KG-GT}
	\label{sec:method2}
	% ═════════════════════════════════════════════════════════════════════════════
	
	\input{method2}
	
	\newpage
	
	% ═════════════════════════════════════════════════════════════════════════════
	\section{Method~3: Phase-Targeted Lightweight Network (PTL-Net)}
	\label{sec:method3}
	% ═════════════════════════════════════════════════════════════════════════════
	
	\input{method3}
	
	\newpage
	
	% ═════════════════════════════════════════════════════════════════════════════
	\section{Method~4: Unified Synergy-Phase-Transformer (USPT)}
	\label{sec:method4}
	% ═════════════════════════════════════════════════════════════════════════════
	
	\input{method4}
	
	\newpage
	
	% ═════════════════════════════════════════════════════════════════════════════
	\section{Comparative Analysis}
	% ═════════════════════════════════════════════════════════════════════════════
	
	\subsection{Baseline and Cross-Method Comparison}
	
	\begin{table}[H]
		\centering
		\caption{Comparison of all proposed methods with baseline approaches.
			Expected performance ranges are based on the reviewed literature and
			architectural analysis.}
		\label{tab:comparison}
		\renewcommand{\arraystretch}{1.3}
		\begin{tabular}{llcccc}
			\toprule
			\textbf{Method} & \textbf{Input} &
			\textbf{Expected $r$} & \textbf{Expected $R^2$} &
			\textbf{Params} & \textbf{Multi-task} \\
			\midrule
			UKF~\cite{sburlea2021} & EEG (delta) &
			$\approx 0.20$ & --- & Minimal & No \\
			CNN-LSTM~\cite{jain2025} & EEG &
			$\approx 0.50$ & $\approx 0.35$ & $\sim$500K & No \\
			\midrule
			\textbf{M1: KG-GT} & EEG + Kin &
			$>0.55$ & $>0.45$ & $\sim$2M & No \\
			\textbf{M2: NNMF-KG-GT} & EEG + Kin + NNMF &
			$>0.60$ & $>0.50$ & $\sim$2M & No \\
			\textbf{M3: PTL-Net} & 6-ch EEG (phase) &
			$>0.65$ & $>0.55$ & $\sim$300K & Yes \\
			\textbf{M4: USPT} & 6-ch EEG + Kin + NNMF &
			$>0.75$ & $>0.65$ & $\sim$1.5M & Yes \\
			\bottomrule
		\end{tabular}
	\end{table}
	
	\subsection{Feature and Complexity Comparison}
	
	\begin{table}[H]
		\centering
		\caption{Architectural feature comparison across all four methods.}
		\label{tab:feature_comp}
		\renewcommand{\arraystretch}{1.3}
		\begin{tabular}{lcccc}
			\toprule
			\textbf{Feature} & \textbf{M1} & \textbf{M2} & \textbf{M3} & \textbf{M4} \\
			\midrule
			EEG-Kinematic CCA alignment  & \checkmark & \checkmark & $\times$   & $\times$ \\
			EEG-EMG CMC synchronisation  & $\times$   & $\times$   & $\times$   & \checkmark \\
			Transformer encoder  & \checkmark & \checkmark & \checkmark & \checkmark \\
			Kinematic-guided GAT & \checkmark & \checkmark & $\times$   & \checkmark \\
			NNMF synergy graph   & $\times$   & \checkmark & $\times$   & \checkmark \\
			Phase targeting      & $\times$   & $\times$   & \checkmark & \checkmark \\
			Channel selection    & $\times$   & $\times$   & \checkmark & \checkmark \\
			Multi-task (Force)   & $\times$   & $\times$   & \checkmark & \checkmark \\
			Weight classifier    & $\times$   & $\times$   & \checkmark & \checkmark \\
			SoftDTW loss         & \checkmark & \checkmark & $\times$   & \checkmark \\
			Phase-weighted loss  & $\times$   & $\times$   & $\times$   & \checkmark \\
			\bottomrule
		\end{tabular}
	\end{table}
	
	% ═════════════════════════════════════════════════════════════════════════════
	\section{Theoretical Justification for Superior Performance}
	% ═════════════════════════════════════════════════════════════════════════════
	
	\subsection{Superiority Over UKF}
	
	\begin{enumerate}[leftmargin=*, itemsep=4pt]
		\item \textbf{Non-parametric expressivity.} The UKF's MVAR state model
		(Eq.~\ref{eq:ukf_state}) is linear in the state history and imposes a
		Gaussian noise structure. Transformer self-attention is a universal
		function approximator over sequences, capable of representing arbitrary
		multi-lag nonlinear mappings between cortical oscillations and multi-muscle
		co-activation.
		
		\item \textbf{No Gaussian assumption.} EEG amplitude distributions near
		movement onset are non-Gaussian. The UKF's statistical guarantees
		degrade; Transformer-based models impose no distributional assumptions on
		intermediate representations.
		
		\item \textbf{Kinematic integration.} The UKF uses EEG as the sole input.
		The WAY-EEG-GAL dataset provides grip aperture, load force, grip force,
		and 3D finger positions that encode the phase and biomechanical state of
		the grasp. The CCA stage and kinematic-guided GAT make full use of this
		information.
		
		\item \textbf{Muscle synergy modelling.} The UKF treats EMG channels
		independently (or through a fixed linear MVAR). The proposed GAT
		explicitly encodes the coordinated relationships between all five muscles,
		consistent with the Discrete Event Sensory Control (DESC) policy
		governing the GAL task.
		
		\item \textbf{NNMF-derived structure (Methods~2, 4).} NNMF provides
		data-driven synergy priors validated by motor neuroscience, replacing
		the UKF's assumption of channel independence.
	\end{enumerate}
	
	\subsection{Superiority Over CNN-LSTM}
	
	\begin{enumerate}[leftmargin=*, itemsep=4pt]
		\item \textbf{Multi-muscle, multi-condition scope.} RawSeq2SeqNet was
		validated on a single biceps brachii channel in a uniplanar task on one
		subject. All four proposed methods target all five muscles across 3,936
		trials with varying weight and friction.
		
		\item \textbf{Kinematic augmentation.} RawSeq2SeqNet uses raw EEG channels
		without task-variable guidance. CCA provides a task-relevant low-dimensional
		projection maximally correlated with kinematic measurements.
		
		\item \textbf{Anatomical priors.} The CNN in RawSeq2SeqNet treats the 19
		EEG channels as an unstructured feature map. The proposed GAT encodes
		known anatomical synergies (e.g., flexor-extensor antagonism; FD--FDI
		co-activation during precision grasp).
		
		\item \textbf{Full-context temporal attention.} The Transformer attends to
		arbitrary lags within the window simultaneously. The LSTM in RawSeq2SeqNet
		processes sequences with exponential temporal decay (vanishing gradient),
		limiting its effective receptive field.
		
		\item \textbf{Phase-aware processing (Methods~3, 4).} By focusing on
		Load and Hold phases and using correlation-based channel selection,
		PTL-Net and USPT eliminate input noise from irrelevant phases and
		channels, improving the signal-to-noise ratio of the prediction.
		
		\item \textbf{Multi-task learning (Methods~3, 4).} Simultaneous prediction
		of EMG, Grip Force, and weight class provides shared-feature regularisation
		and demonstrates clinical applicability beyond single-output regression.
	\end{enumerate}
	
	% ═════════════════════════════════════════════════════════════════════════════
	\section{Expected Results}
	% ═════════════════════════════════════════════════════════════════════════════
	
	Based on the architectural design and prior literature, we state the following
	empirical hypotheses:
	
	\begin{enumerate}[label=\textbf{H\arabic*.}, leftmargin=*, itemsep=4pt]
		\item Method~4 (USPT) will achieve the highest overall Pearson $r > 0.75$
		across all five EMG channels and 12 subjects, followed by Method~3
		($r > 0.65$), Method~2 ($r > 0.60$), and Method~1 ($r > 0.55$).
		\item Prediction accuracy will be highest for the flexor digitorum and
		first dorsal interosseus (the primary precision grasp muscles),
		consistent with their strong kinematic correlation during the Load phase.
		\item Phase-specific analysis will show significantly higher $r$ and $R^2$
		during Load and Hold phases compared to Reach and Replace, validating
		the phase-targeting strategy of Methods~3 and~4.
		\item The NNMF-derived synergy graph (Methods~2, 4) will produce
		consistently higher $r$ than the anatomically fixed graph (Method~1),
		with synergy similarity index $\mathrm{SSI} > 0.80$ across subjects.
		\item The weight classifier (Methods~3, 4) will achieve $> 85\%$
		accuracy and $F_1 > 0.83$ across three weight classes, demonstrating
		the feasibility of EEG-based adaptive prosthetic grip control.
		\item The SoftDTW component will reduce DTW temporal distance by at least
		$15\%$ relative to MSE-only training, without degrading nRMSE.
		\item Method~3 (PTL-Net) will have the lowest computational cost
		($< 30\%$ of Method~1 FLOPs) while maintaining competitive accuracy,
		demonstrating the efficiency of phase-targeted channel selection.
	\end{enumerate}
	
	% ═════════════════════════════════════════════════════════════════════════════
	\section{Discussion}
	% ═════════════════════════════════════════════════════════════════════════════
	
	The four-method comparative framework presented here systematically
	investigates the contribution of each architectural innovation to
	EEG-to-EMG prediction quality.
	
	\paragraph{Method~1 (KG-GT)} establishes the baseline deep-learning
	architecture with CCA alignment, Transformer encoding, and
	kinematic-conditioned GAT. The kinematic conditioning of GAT attention
	is a theoretical contribution beyond existing graph-based EMG predictors:
	$\alpha_{ij}^t$ is recomputed at every time step from $\bm{k}_t$, making
	inter-muscle synergy representation dynamic and context-dependent.
	
	\paragraph{Method~2 (NNMF-KG-GT)} replaces anatomical graph assumptions
	with data-driven NNMF synergy weights, connecting decades of motor
	neuroscience research~\cite{tresch2006, pierella2020} to modern
	graph neural networks. The cosine similarity of NNMF loadings provides
	a principled initialisation that biases the GAT towards physiologically
	meaningful muscle coordination patterns.
	
	\paragraph{Method~3 (PTL-Net)} demonstrates that simpler, targeted
	approaches can be highly effective. By exploiting phase-specific
	EEG--EMG coupling and reducing input dimensionality through channel
	selection, PTL-Net achieves competitive accuracy with dramatically lower
	computational cost. The multi-task architecture (EMG + Force + weight
	class) directly addresses prosthetic application requirements.
	
	\paragraph{Method~4 (USPT)} combines all innovations into a unified
	framework, with phase embeddings, NNMF-initialised GAT, and
	phase-weighted loss. The staged training protocol ensures stable
	convergence of this complex architecture.
	
	The perturbation trials of WAY-EEG-GAL provide a unique test of motor
	adaptation decoding. EEG channels Pz and C4 are known from the original
	paper to show significant alpha and beta power changes approximately 200~ms
	after unexpected weight changes~\cite{luciw2014}. If USPT achieves
	measurably different prediction accuracy on these trials relative to stable
	trials, it would indicate that the architecture is capturing the EEG correlates
	of sensorimotor adaptation---a clinically significant capability for prosthetic
	devices that must respond to objects with varying properties.
	
	\paragraph{Limitations.}
	(1)~CCA assumes a stationary linear EEG--kinematic relationship; adaptive CCA
	or online recalibration should be explored for non-stationary real-world use.
	(2)~The graph structure is currently defined either anatomically or via NNMF;
	end-to-end graph learning may discover unexpected synergistic relationships.
	(3)~All experiments are offline; real-time implementation introduces latency
	constraints on window length and model complexity.
	(4)~The recommended channel set (C3, C4, Cz, CP3, CP4, Fz) is based on
	general sensorimotor topography; subject-specific channel optimisation may
	further improve individual performance.
	
	% ═════════════════════════════════════════════════════════════════════════════
	\section{Conclusion}
	% ═════════════════════════════════════════════════════════════════════════════
	
	This article proposes a multi-method comparative framework for continuous
	prediction of EMG envelopes from scalp EEG and kinematic data during
	grasp-and-lift tasks on the WAY-EEG-GAL dataset. Four architectures of
	increasing complexity address eight identified research gaps:
	
	\textbf{Method~1 (KG-GT)} combines CCA-based EEG--kinematic alignment,
	Transformer-based nonlinear temporal encoding, kinematic-conditioned dynamic
	GAT, and hybrid MSE--SoftDTW training loss.
	
	\textbf{Method~2 (NNMF-KG-GT)} enhances Method~1 with NNMF-derived muscle
	synergy weights for data-driven, biologically interpretable graph
	initialisation.
	
	\textbf{Method~3 (PTL-Net)} achieves competitive performance with reduced
	computational cost by targeting Load and Hold phases, selecting 6 optimal
	EEG channels, and performing multi-task prediction of EMG, Grip Force, and
	object weight class.
	
	\textbf{Method~4 (USPT)} unifies all innovations into a single framework
	with phase-aware Transformer encoding, NNMF-initialised kinematic-guided
	GAT, triple-head decoding, and phase-weighted loss.
	
	Theoretical arguments and prior empirical results support the expectation
	of substantially improved accuracy over both UKF-based state estimators and
	conventional CNN-LSTM architectures across multi-muscle, multi-condition
	grasping tasks. The weight classification capability of Methods~3 and~4
	directly demonstrates clinical applicability for adaptive prosthetic control.
	
	% ═════════════════════════════════════════════════════════════════════════════
	\bibliographystyle{ieeetr}
	\begin{thebibliography}{99}
		
		\bibitem{luciw2014}
		M.~D. Luciw, E.~Jarocka, and B.~B. Edin,
		``Multi-channel EEG recordings during 3,936 grasp and lift trials with
		varying weight and friction,''
		\textit{Scientific Data}, vol.~1, p.~140047, 2014.
		
		\bibitem{nakagome2017}
		S.~Nakagome, T.~P. Luu, J.~A. Brantley, and J.~L. Contreras-Vidal,
		``Prediction of EMG envelopes of multiple terrains over-ground walking
		from EEG signals using an Unscented Kalman Filter,''
		in \textit{Proc.\ IEEE Int.\ Conf.\ Syst., Man, Cybern.\ (SMC)}, 2017.
		
		\bibitem{sburlea2021}
		A.~I. Sburlea, N.~Butturini, and G.~R. M\"{u}ller-Putz,
		``Predicting EMG envelopes of grasping movements from EEG recordings
		using Unscented Kalman Filtering,''
		in \textit{Proc.\ Annu.\ Meeting Austrian Soc.\ Biomed.\ Eng.}, 2021.
		
		\bibitem{jain2025}
		P.~Jain, V.~Bharadiya, U.~Chittimalla, Y.~Anand, and M.~Rao,
		``EEG-Based Surface EMG Reconstruction Using Deep Sequence Learning for
		Upper Limb Motor Activity,'' 2025.
		
		\bibitem{li2025}
		M.~Li, Z.~Wei, Z.-Q. Zhang, S.~Ma, and S.~Q. Xie,
		``Continuous Joint Kinematics Prediction Using GAT-LSTM Framework Based
		on Muscle Synergy and Sparse sEMG,''
		\textit{IEEE Trans.\ Neural Syst.\ Rehabil.\ Eng.}, vol.~33, pp.~1763--1773,
		2025.
		
		\bibitem{cuturi2017}
		M.~Cuturi and M.~Blondel,
		``Soft-DTW: A differentiable loss function for time-series,''
		in \textit{Proc.\ 34th Int.\ Conf.\ Mach.\ Learn.\ (ICML)},
		vol.~70, pp.~894--903, 2017.
		
		\bibitem{vaswani2017}
		A.~Vaswani et al.,
		``Attention is all you need,''
		in \textit{Adv.\ Neural Inf.\ Process.\ Syst.\ (NeurIPS)}, vol.~30, 2017.
		
		\bibitem{velickovic2018}
		P.~Veli\v{c}kovi\'{c} et al.,
		``Graph Attention Networks,''
		in \textit{Proc.\ Int.\ Conf.\ Learn.\ Represent.\ (ICLR)}, 2018.
		
		\bibitem{pierella2020}
		C.~Pierella, E.~Pirondini, N.~Kinany, M.~Coscia, C.~Giang,
		J.~Miehlbradt, C.~Magnin, P.~Nicolo, S.~Dalise, G.~Sgherri,
		C.~Chisari, D.~Van~De~Ville, A.~Guggisberg, and S.~Micera,
		``A multimodal approach to capture post-stroke temporal dynamics
		of recovery,''
		\textit{J.\ Neural Eng.}, vol.~17, no.~4, p.~045002, 2020.
		
		\bibitem{lee1999}
		D.~D. Lee and H.~S. Seung,
		``Learning the parts of objects by non-negative matrix factorization,''
		\textit{Nature}, vol.~401, no.~6755, pp.~788--791, 1999.
		
		\bibitem{tresch2006}
		M.~C. Tresch, V.~C.~K. Cheung, and A.~d'Avella,
		``Matrix factorization algorithms for the identification of muscle
		synergies: Evaluation on simulated and experimental data sets,''
		\textit{J.\ Neurophysiol.}, vol.~95, no.~4, pp.~2199--2212, 2006.
		
		\bibitem{cmc_paper}
		``Coherence-based connectivity analysis of EEG and EMG signals during reach-to-grasp movement involving two weights.''
		
		\bibitem{esi_gal_2025}
		``Various time-lagged and window sizes are analyzed for hand kinematics prediction: ESI-GAL EEG Source Imaging-based Kinematics Parameter Estimation for Grasp and Lift Task,'' 2025.
		
	\end{thebibliography}
	
\end{document}